%! AP Statistics — Unit 9, Lesson 9.3 — Group Activity
\documentclass[10pt]{article}
\usepackage{apstats-article}
\usepackage{apstats-boxes}

\begin{document}

\pageheader{Unit 9, Lesson 9.3}{Group Activity}
\namepartnerperiod

\vspace{0.1in}
Three research teams each fit a regression model and reported a 95\% confidence interval
for the slope. The intervals are already computed --- your job is to interpret them and
evaluate researcher claims.

\vspace{0.08in}
\begin{center}
\renewcommand{\arraystretch}{1.4}
\begin{tabular}{clllcc}
    \textbf{Study} & \textbf{Explanatory ($x$)} & \textbf{Response ($y$)} &
    \textbf{Population} & \textbf{95\% CI} & \textbf{$n$} \\
    \hline
    A & Daily screen time (hr) & Sleep duration (hr) & High school students & $(-0.31,\ -0.05)$ & 45 \\
    B & Advertising spend (\$1000s) & Monthly sales (\$1000s) & Small retailers & $(-0.8,\ 2.1)$ & 28 \\
    C & Elevation (100 ft) & Avg.\ July temperature ($^\circ$F) & Weather stations & $(-3.2,\ -0.9)$ & 20 \\
\end{tabular}
\end{center}

\vspace{0.1in}
\textbf{Directions:} Complete the tier that matches your current level of understanding.

\begin{tcolorbox}[colback=white, colframe=black!40,
    title=\textbf{Tier R --- Remediate},
    fonttitle=\bfseries, arc=2mm, left=3mm, right=3mm, top=2mm, bottom=2mm]

Use \textbf{Study A} only: $x =$ daily screen time, $y =$ sleep duration,
population = high school students, CI $= (-0.31,\ -0.05)$.

\textbf{1.} Interpret the CI. Fill in the sentence frame:

\vspace{0.05in}
``We are 95\% confident that the true slope relating \blank{4.5cm} to \blank{4.5cm}
for \blank{4cm} is between \blank{2cm} and \blank{2cm} \blank{3cm}.''

\vspace{0.1in}
\textbf{2.} Is 0 in the interval $(-0.31,\ -0.05)$? \quad \textbf{Yes} \quad \textbf{No}
\quad (circle one)

\vspace{0.05in}
Therefore, there \blank{2cm} evidence that $\beta \ne 0$.

\vspace{0.1in}
\textbf{3.} A researcher claims $\beta = -0.1$. Is $-0.1$ in the interval $(-0.31,\ -0.05)$?

\quad \textbf{Yes} \quad \textbf{No} \quad (circle one)

\vspace{0.05in}
Is the claim $\beta = -0.1$ consistent with the data? \blank{5cm}

\end{tcolorbox}

\begin{tcolorbox}[colback=white, colframe=black!40,
    title=\textbf{Tier A --- Approaching Proficiency},
    fonttitle=\bfseries, arc=2mm, left=3mm, right=3mm, top=2mm, bottom=2mm]

For each of the three studies (A, B, and C):

\begin{enumerate}[label=\alph*.]

  \item Write a complete interpretation of the 95\% CI for slope in context.

        Study A: \writelines{2}
        Study B: \writelines{2}
        Study C: \writelines{2}

  \item Does the interval provide evidence that $\beta \ne 0$? Justify briefly.

        Study A: \writelines{2}
        Study B: \writelines{2}
        Study C: \writelines{2}

  \item Evaluate the researcher's claim for each study.
        (Study A: $\beta = 0$; \quad Study B: $\beta = 1.5$; \quad Study C: $\beta = -2.0$)

        Study A: \writelines{2}
        Study B: \writelines{2}
        Study C: \writelines{2}

\end{enumerate}
\end{tcolorbox}

\begin{tcolorbox}[colback=white, colframe=black!40,
    title=\textbf{Tier E --- Extension},
    fonttitle=\bfseries, arc=2mm, left=3mm, right=3mm, top=2mm, bottom=2mm]

Complete all of Tier A above. Then answer the following extension questions.

\textbf{Bonus 1.} Study B's 95\% CI is $(-0.8,\ 2.1)$. Write a one-sentence explanation
of why this interval does \textbf{not} provide evidence of a linear relationship between
advertising spend and monthly sales.

\writelines{2}

\textbf{Bonus 2.} Study C used $n = 20$ weather stations. If the researchers repeated
the study with $n = 80$ stations (same sample standard deviation and confidence level),
would the new CI likely contain $-2.0$? Would it be wider or narrower than $(-3.2,\ -0.9)$?
Explain using what you know about sample size and CI width.

\writelines{3}

\end{tcolorbox}

\end{document}
